\chapter{Ordination und Einführungen}
\quellseite{10}

Im Auftrag des Ehrw. Missionsdirectors P. Mahler wurde Cand. Emil Schulz am 17. Sonnt. n. Trin. in seiner Gemeinde zu Sao Pedro, Brasilien, ordinirt und eingeführt von H. Stieme.

Im Auftrag des Ehrw. Präses Becker wurde P. Friedrich Schwarz am 3. Sonnt. n. Epiph. in seiner Gemeinde zu Carroll, Nebr., eingeführt von C. Gutknecht.

\quellseite{11}

Im Auftrag des Ehrw. Präses Runkel wurde P. M. Liebe am Sonnt. Sexagesimä in seiner Gemeinde im Richmond-District, von San Francisco, Cal., unter Assistenz der PP. Schröder, Block und Thiede eingeführt von G. A. Berntahl.

Im Auftrag des Ehrw. Präses C. Hafner wurde P. Otto Winnat am Sonnt. Quinquagesimä in seiner Gemeinde bei Bern, Kans., eingeführt von P. W. Graf.

Im Auftrag des Ehrw. Präses Wegener wurde P. C. Buchschacher am Sonnt. Quinquagesimä in seiner Gemeinde zu Swiss Alp, Tex., unter Assistenz P. C. Buchschachers sen. eingeführt von R. Osthoff.

Im Auftrag des Ehrw. Präses Seuel wurde P. C. D. Griese am Sonnt. Synodcal in seinen Gemeinden zu Polar und Norwood, Wis., unter Assistenz P. D. Lichts eingeführt von D. Grimm.

Am Sonnt. Sexagesimä wurde Lehrer B. M. Hemmann als Lehrer der zweiten Schulklasse der Gemeinde zu Altenburg, Mo., eingeführt von A. A. Schormann.
